\subsection{Cobertura y validez de la campaña}
\label{subsec:resultados-cobertura}

% Redactar: abrir el capítulo explicando que la comparación se apoya en una
% campaña cerrada y trazable, no en ejemplos aislados. El tono debe ser de
% auditoría experimental: qué entra, qué queda fuera y por qué.
% Evidencia: comparison_all_runs.csv contiene 264 corridas comparables:
% 184 en test (23 escenarios base x 8 controladores) y 80 en OOD
% (10 escenarios base x 8 controladores). neural_position queda fuera por no
% disponer de evidencia consolidada.

La campaña comparativa queda resumida en la \autoref{tab:cobertura-campana}.
Los artefactos completos de generación, consolidación y visualización se
documentan en el \apartadoref{anejo:comandos}.

\begin{table}[H]
\caption{Cobertura de la campaña comparativa principal.}
\label{tab:cobertura-campana}
\begin{tabular}{lrrrrr}
\toprule
Split & Escenarios base & Controladores & Corridas & Éxito misión [\%] & Seguridad [\%] \\
\midrule
\texttt{test} & 23 & 8 & 184 & 99,5 & 99,5 \\
\texttt{ood} & 10 & 8 & 80 & 66,3 & 76,3 \\
\bottomrule
\end{tabular}
\end{table}

% Redactar: mencionar que no se eliminan episodios fallidos antes de agregar.
% Dato útil: en test solo falla circle_g04_P5_combined_s1083 con PD Lissajous
% transferido por límite de actitud. En OOD los fallos se concentran en
% lemniscatas extremas, compuestos agresivos y helix_tight_aggressive.
% Nota: no extenderse aquí; las causas se retoman con las figuras OOD.

\subsection{Referencia clásica y transferencia entre familias}
\label{subsec:resultados-referencia-clasica}

% Redactar: presentar esta parte como la referencia de ingeniería. Antes de
% juzgar redes, se muestra qué ocurre cuando un PD especializado se transfiere
% sin reajuste a otras familias.
% Lectura esperada: la diagonal no domina en todas las familias. Hay
% especialización parcial, no confirmación total. El PD Lissajous transfiere
% especialmente bien en hold y waypoint.

\begin{figure}[H]
\centering
\includegraphics[width=0.72\textwidth]{res_pid_transfer_matrix.pdf}
\caption{Transferencia cruzada de los controladores PD en el split
\texttt{test}. Las filas indican la familia usada para la sintonización y las
columnas la familia evaluada. La diagonal marca el uso nativo de cada
controlador; las excepciones fuera de la diagonal deben interpretarse como
transferencia favorable, no como reajuste posterior.}
\label{fig:res-pid-transfer-matrix}
\end{figure}

% Redactar: usar la figura para defender una idea matizada. No decir que el PD
% nativo es siempre óptimo. Datos útiles: circle nativo 0,094 m frente a 0,162 m
% con PD Hold; lissajous nativo 0,071 m; hold mejora con PD Lissajous 0,044 m
% frente a 0,057 m; waypoint mejora con PD Lissajous 0,048 m frente a 0,066 m.

\subsection{Aprendizaje neuronal en condiciones conocidas}
\label{subsec:resultados-neuronal-id}

% Redactar: conectar entrenamiento supervisado y bucle cerrado. La tabla no
% demuestra por sí sola buen control, pero permite comparar fidelidad de fuerza
% con el comportamiento posterior.

\begin{table}[H]
\caption{Fidelidad supervisada de fuerza y señales de protección asociadas.}
\label{tab:fidelidad-supervisada}
\begin{tabular}{lrrrr}
\toprule
Arquitectura & MSE norm. & RMSE fuerza norm. [N] & Clip incl. previo [\%] & Clip incl. cerrado [\%] \\
\midrule
MLP  & 0,0185 & 0,083 & 0,057 & 0,062 \\
GRU  & 0,0102 & 0,053 & 0,000 & 0,000 \\
LSTM & 0,0111 & 0,055 & 0,000 & 0,000 \\
\bottomrule
\end{tabular}
\end{table}

% Redactar: explicar que la fidelidad supervisada no se traduce de forma
% automática en menor RMSE en bucle cerrado. Esta es una de las ideas centrales.

\begin{figure}[H]
\centering
\includegraphics[width=0.92\textwidth]{res_id_rmse_family.pdf}
\caption{RMSE medio en condiciones conocidas para el PD representativo y las
tres arquitecturas neuronales outer-force. Se muestra únicamente el seguimiento,
ya que el éxito de misión en \texttt{test} apenas discrimina entre
controladores.}
\label{fig:res-id-rmse-family}
\end{figure}

% Redactar: destacar que MLP es competitivo en hold/lissajous, GRU mejora en
% circle frente a LSTM, y LSTM no domina pese a su fidelidad supervisada. Evitar
% frases de superioridad global.

\subsection{Generalización fuera de distribución}
\label{subsec:resultados-transferencia-ood}

% Redactar: esta debe ser la parte visualmente central del capítulo. La figura
% por familia da la lectura rápida; la matriz por escenario demuestra que no se
% esconden casos extremos ni heterogeneidad interna.

\begin{figure}[H]
\centering
\includegraphics[width=0.92\textwidth]{res_ood_rmse_family.pdf}
\caption{RMSE medio fuera de distribución por familia de escenario. La escala se
mantiene lineal para conservar la interpretación física en metros; las
lemniscatas y lissajous extremas elevan el error de algunas arquitecturas y
deben leerse junto con el desglose por escenario.}
\label{fig:res-ood-rmse-family}
\end{figure}

\begin{figure}[H]
\centering
\includegraphics[width=0.86\textwidth]{res_ood_scenario_matrix.pdf}
\caption{Desglose OOD por escenario y controlador para las cuatro referencias
principales. El número de cada celda es el RMSE en metros; el asterisco indica
que la misión no se completó aunque el episodio se conserva en la agregación.}
\label{fig:res-ood-scenario-matrix}
\end{figure}

% Redactar: separar claramente tres fenómenos: OOD moderado donde las redes son
% competitivas; geometrías extremas donde aparecen errores grandes; y fallos de
% misión que no deben ocultarse en medias. La matriz permite citar casos
% concretos sin cargar el texto con todos los números.

\begin{figure}[H]
\centering
\includegraphics[width=0.76\textwidth]{res_ood_termination_summary.pdf}
\caption{Tipos de terminación en la batería OOD para los controladores
principales. Esta figura complementa el RMSE: una trayectoria puede tener error
moderado y, aun así, no satisfacer el criterio de misión.}
\label{fig:res-ood-termination-summary}
\end{figure}

% Redactar: usar esta figura para explicar por qué no basta con una media de
% error. Distinguir tiempo, saturación y actitud como modos de fallo con
% implicaciones físicas distintas.

\subsection{Caso representativo de arquitectura}
\label{subsec:resultados-caso-representativo}

% Redactar: esta figura no sustituye a la estadística agregada. Sirve para que
% el lector vea el fenómeno dinámico: dos arquitecturas con la misma interfaz
% outer-force pueden producir seguimiento muy distinto.

\begin{figure}[H]
\centering
\includegraphics[width=0.88\textwidth]{res_trajectory_lemniscate_mlp_lstm.pdf}
\caption{Escenario OOD \texttt{lemniscate\_fast\_center\_yaw}: comparación
visual entre MLP y LSTM. El panel superior muestra el plano horizontal y el
inferior el error de posición a lo largo del episodio.}
\label{fig:res-trajectory-lemniscate-mlp-lstm}
\end{figure}

% Redactar: defender aquí la idea "la arquitectura decide" sin convertirla en
% regla universal. La figura ilustra un caso fuerte; la evidencia general está
% en las figuras agregadas y en el atlas del anejo.

\subsection{Actuación, protecciones y límites compartidos}
\label{subsec:resultados-protecciones}

% Redactar: cerrar la parte experimental recordando que la seguridad no procede
% solo de la red. Hay mixer, saturación, límites de fuerza y degradación experta.
% No presentar clipping frecuente como éxito del aprendizaje.

\begin{figure}[H]
\centering
\includegraphics[width=0.78\textwidth]{res_protections_ood.pdf}
\caption{Activación media de protecciones y degradación experta en OOD. La
figura ayuda a separar error de seguimiento, demanda de control y actuación de
los límites compartidos del sistema.}
\label{fig:res-protections-ood}
\end{figure}

% Redactar: explicar que MLP concentra clipping de norma/inclinación en OOD,
% mientras que la degradación experta aparece también en referencias clásicas.
% Usar esta figura como transición natural hacia la discusión.

\subsection{Lecturas principales de la campaña}
\label{subsec:resultados-sintesis}

% Redactar: esta tabla es un puente hacia discusión. Mantenerla como síntesis
% sobria: evidencia observada + lectura prudente. Evitar convertirla en una
% checklist explícita de H1--H5, aunque cubra esas preguntas por debajo.

\begin{table}[H]
\caption{Lecturas principales que debe desarrollar la redacción final.}
\label{tab:sintesis-resultados}
\begin{tabularx}{\textwidth}{p{0.27\textwidth} X X}
\toprule
Eje de lectura & Evidencia observada & Lectura acotada \\
\midrule
Especialización clásica & La matriz PD muestra diagonal favorable en parte de
las familias y excepciones claras en \texttt{hold} y \texttt{waypoint}. &
La especialización existe, pero la transferencia entre familias no es
uniformemente mala. \\
Imitación en \texttt{test} & MLP, GRU y LSTM completan las misiones conocidas,
con RMSE dependiente de familia. & La fidelidad supervisada no basta para
predecir el comportamiento en bucle cerrado. \\
Generalización OOD & Los agregados y la matriz por escenario muestran fallos y
errores altos en geometrías extremas. & La política común es competitiva en
casos moderados, no superior de forma global. \\
Memoria temporal & GRU y LSTM no dominan sistemáticamente a MLP ni entre sí. &
La memoria recurrente introduce una sensibilidad arquitectónica que debe
interpretarse por familia y escenario. \\
Protecciones & Clipping, saturación y degradación aparecen de forma desigual
entre controladores y escenarios. & La seguridad observada pertenece al sistema
híbrido completo, no únicamente a la red. \\
\bottomrule
\end{tabularx}
\end{table}

% Redactar: al cerrar el capítulo, remitir al anejo como atlas visual. Las
% figuras de trayectorias del anejo no son demostración comparativa, sino una
% muestra de la variedad de escenarios y del trabajo experimental realizado.
% Referencia útil: \apartadoref{anejo:comandos}.
